%==============================================================
%  CHAPTER 7 - REMOTE ACCESS      (skeleton: fill in and delete the TODO boxes)
%  Pattern: copy any \shot line, change the file name, caption and label.
%  Screenshots go in images/ using the file name written here (remote_01.png, ...).
%==============================================================
\chapter{Remote Access Configuration}
\label{chap:remote}

\section{What is Remote Access?}
\label{sec:remote-what}
\todo{Write 3--5 sentences defining the Remote Access role (and VPN / DirectAccess, whichever was configured) and cite a source with \texttt{\textbackslash cite\{key\}}.}

\section{Why is Remote Access Important?}
\label{sec:remote-why}
\todo{Explain why the organisation needs it (secure connectivity for remote users) and how it fits with AD DS and Group Policy.}

\section{Security Risks if Not Set Up or Set Up Incorrectly}
\label{sec:remote-risks}
\todo{Describe the risks of leaving the role out or misconfiguring it.}

\section{Step-by-Step Configuration Guide}
\label{sec:remote-guide}
\todo{Add a short table of the key values used (see Table~\ref{tab:config} in Chapter~\ref{chap:adds} for the format), then one subsection per step. Each step: one paragraph of text that refers to its screenshot with \texttt{\textbackslash autoref\{fig:remote-01\}}, followed by the \texttt{\textbackslash shot} line.}

\subsection{Step 1: Install the Remote Access Role}
\label{sec:remote-step1}
\todo{Describe the step here.}

\shot{remote_01.png}{Replace this caption (Remote Access, step 1)}{fig:remote-01}
\shot{remote_02.png}{Replace this caption (Remote Access, step 1 continued)}{fig:remote-02}

\subsection{Step 2: Configure the Service}
\label{sec:remote-step2}
\todo{Describe the step here.}

\shot{remote_03.png}{Replace this caption (Remote Access, step 2)}{fig:remote-03}

\section{Testing}
\label{sec:remote-testing}
\todo{Explain how the service was tested and what the result was.}

\shot{remote_test_01.png}{Replace this caption (Remote Access test result)}{fig:remote-test-01}

\section{Troubleshooting: Possible Sticking Points \& Fixes}
\label{sec:remote-troubleshooting}
\issue{1. Replace with the problem title}
{Describe what the user sees.}
{Describe how to fix it.}
